\documentclass[11pt,a4paper]{article}

\usepackage[margin=0.85in]{geometry}
\usepackage[T1]{fontenc}
\usepackage{lmodern}
\usepackage{microtype}
\usepackage{graphicx}
\usepackage{booktabs}
\usepackage{tabularx}
\usepackage{enumitem}
\usepackage{xcolor}
\usepackage{hyperref}
\usepackage{amsmath}
\usepackage{array}
\usepackage{titlesec}

\definecolor{IITKBlue}{HTML}{17365D}
\definecolor{Muted}{HTML}{555555}
\hypersetup{
  colorlinks=true,
  linkcolor=IITKBlue,
  urlcolor=IITKBlue,
  citecolor=IITKBlue,
  pdftitle={Video Transcription with Confidence Analytics},
  pdfsubject={CS698B Fundamentals of Data Engineering-I Final Project Proposal}
}
\titleformat{\section}{\Large\bfseries\color{IITKBlue}}{\thesection.}{0.55em}{}
\titleformat{\subsection}{\large\bfseries\color{IITKBlue}}{\thesubsection}{0.55em}{}
\setlist[itemize]{leftmargin=1.4em,itemsep=2pt,topsep=3pt}
\setlist[enumerate]{leftmargin=1.6em,itemsep=2pt,topsep=3pt}
\renewcommand{\arraystretch}{1.18}

\title{\vspace{-1.1cm}\textbf{\Huge Video Transcription with Confidence Analytics}\\[0.35em]
\large CS698B -- Fundamentals of Data Engineering-I\\
\normalsize Final Project Proposal}
\author{}
\date{}

\begin{document}
\maketitle
\vspace{-1.2em}

\begin{abstract}
We propose a web application that converts uploaded, pre-recorded videos into
speaker-attributed, time-aligned transcripts and communicates how reliable
different parts of the generated transcript are. In addition to transcription,
the system estimates word- and segment-level confidence, evaluates those
estimates against reference transcripts on a benchmark collection, and
visualizes uncertain regions so that users can prioritize manual review.
The project combines an audio ETL pipeline, pretrained speech models,
confidence estimation and calibration, an interactive visualization
dashboard, persistent application services, and reproducible software
engineering practices.
\end{abstract}

\section{Introduction and Motivation}

Conventional speech-to-text systems primarily return a transcript. A transcript
alone, however, does not tell a user which words are likely to be correct,
which regions are affected by noise or overlapping speech, or where manual
verification is most important. This limitation is particularly relevant to
recordings containing multiple speakers, variable acoustic conditions, and
domain-specific vocabulary.

Our project, \textbf{Video Transcription with Confidence Analytics}, addresses
this problem through a web application for pre-recorded video and audio.
The application extracts the audio, processes it through a speech-intelligence
pipeline, and returns a time-aligned transcript with speaker attribution and
confidence information. The dashboard highlights uncertain words and
segments, provides audio-linked inspection, and summarizes the overall
transcription quality.

The project uses publicly available recordings with manually prepared
reference transcripts for development and quantitative evaluation. These
references are available for the benchmark collection, not for arbitrary
user uploads. Therefore, uploaded recordings are transcribed without assuming
that a reference transcript exists. Confidence estimation for user uploads
uses signals available from the model and audio; its calibration and measured
reliability are assessed separately on the labelled benchmark data.

The scope is limited to \textbf{pre-recorded media}. Live transcription is
not part of the current project.

\section{Objectives}

The main objectives are:
\begin{enumerate}
  \item Build a reproducible ETL pipeline that converts uploaded media into
  validated, standardized audio and structured transcript records.
  \item Generate accurate word-level transcripts, timestamps, and speaker
  attribution using pretrained speech models.
  \item Estimate confidence at word and segment level, and study whether
  these confidence values correspond to observed transcription correctness.
  \item Evaluate transcription and confidence quality using reference data and
  appropriate metrics.
  \item Develop an accessible dashboard for exploring transcripts, audio
  timelines, speakers, confidence trends, and evaluation results.
  \item Deliver the system as a usable application with persistent records,
  documented interfaces, automated tests, and reproducible configuration.
\end{enumerate}

\section{Project Components}

\subsection{Data Component -- ETLS}

The data component manages media ingestion, transformation, validation,
structured storage, and reproducibility.

\paragraph{Ingestion.}
Users upload a pre-recorded video or audio file. The system validates the
file type and media stream, extracts the audio track from video, and records
metadata such as filename, duration, sample rate, channel count, and file
hash. Supported input formats are intended to include common formats such as
MP4, MOV, MP3, WAV, M4A, and FLAC, subject to the installed media decoder.

\paragraph{Transformation.}
The audio is converted to a consistent representation (16\,kHz, mono PCM WAV)
for downstream processing. The pipeline performs audio validation and
loudness normalization. Voice Activity Detection (VAD) identifies speech
regions and supports segmentation of long recordings. The system retains
timestamps so that transcript words and speaker turns can be mapped back to
the original media timeline. Noise reduction may be evaluated as an optional
preprocessing experiment; it is not assumed to improve recognition in every
recording condition.

\paragraph{Loading and storage.}
The application stores recording metadata, processing-job status, word-level
transcription results, speaker turns, confidence values, and user revisions
in a structured database. SQLite is suitable for local development, with
PostgreSQL as an alternative for multi-user deployment. Intermediate
tabular evaluation data may be stored in Parquet or CSV. Original uploads
and generated media artifacts are stored separately from relational metadata.


\subsection{Adaptive / AI--ML Component}

This component performs speech recognition, speaker analysis, and
confidence estimation. It uses pretrained models rather than training a
large speech model from scratch.

\paragraph{Speech processing pipeline.}
The proposed processing sequence is:
\begin{enumerate}
  \item \textbf{Audio preprocessing:} standardize the audio format and
  loudness while retaining the original timeline.
  \item \textbf{Voice Activity Detection:} use a high-recall VAD model to
  identify speech regions and split long regions into manageable chunks.
  \item \textbf{Automatic Speech Recognition:} use Faster-Whisper
  (Large-v3, subject to available hardware) to produce the transcript.
  \item \textbf{Forced alignment:} use a Wav2Vec2-based aligner through
  WhisperX to estimate word-level start and end times.
  \item \textbf{Speaker diarization:} use NeMo MSDD with TitaNet speaker
  representations, conditioned on speech regions, to estimate who spoke
  when.
  \item \textbf{Speaker--word fusion:} associate each aligned word with the
  diarization timeline using temporal overlap and a word-midpoint rule.
  Consecutive words assigned to the same speaker are grouped into turns.
\end{enumerate}

The pipeline is modular: individual stages can be evaluated and replaced
without changing the data model or visualization layer. The selected models
are pretrained; the project focuses on integration, evaluation, confidence
modelling, and application-level intelligence.

\paragraph{Confidence estimation.}
A model's raw decoding score is not automatically a calibrated probability
that a word is correct. We will therefore distinguish between:
\begin{itemize}
  \item \textbf{Raw confidence features:} available token/word likelihood
  information and model scores, where exposed by the ASR implementation;
  audio-quality indicators; speech-region characteristics; and indicators
  of difficult conditions such as overlap or unusually short utterances.
  \item \textbf{Calibrated confidence:} a score mapped to the empirical
  probability of correctness using labelled benchmark examples.
\end{itemize}

Word correctness will be derived by aligning the system hypothesis with the
reference transcript. A calibration model will be fitted using benchmark
predictions and evaluated on held-out recordings or recording-level
cross-validation. The project will investigate a suitable calibration method,
such as temperature scaling or isotonic regression, depending on the
available sample size and score distribution. Calibration must be evaluated
on data not used to fit the calibration mapping.

Confidence will be aggregated to segment level using a documented rule,
such as the mean or a robust summary of the constituent word scores. The
dashboard will label these values as estimates and will not present an
uncalibrated model score as a guaranteed probability of correctness.

\paragraph{Adaptive feedback.}
Users may correct words, speaker labels, or turn boundaries. The application
records the original model output and each subsequent revision. Corrections
can be analysed to identify recurring lexical errors and domain-specific
terms. Where supported by the ASR model, verified vocabulary can be supplied
as a contextual prompt in later runs. The feedback mechanism is initially
human-in-the-loop: corrections are collected and analysed; automatic model
retraining is not assumed.

\subsection{Visualization Component -- VIZ}

The visualization component makes the transcript and its reliability
interpretable and actionable.

\begin{itemize}
  \item \textbf{Transcript view:} display the transcript by speaker and
  time range, with word-level confidence markers.
  \item \textbf{Confidence highlighting:} distinguish low-, medium-, and
  high-confidence words using colour together with symbols, labels, or
  underline patterns. Colour is never the only indicator.
  \item \textbf{Audio timeline:} display the waveform or speech timeline,
  speaker-turn regions, and confidence information against media time.
  \item \textbf{Synchronized playback:} allow a user to select a word or
  segment, play the corresponding audio region, and follow the active word
  during playback.
  \item \textbf{Summary analytics:} report recording duration, word count,
  number of detected speakers, average confidence, and the number and
  proportion of low-confidence words.
  \item \textbf{Evaluation analytics:} show WER and calibration results,
  including reliability diagrams, confidence distributions, and breakdowns
  by recording or acoustic condition where the available labels support them.
  \item \textbf{Review and export:} allow users to inspect uncertain regions,
  make corrections, and export transcript data and subtitles.
\end{itemize}

The dashboard will use a high-contrast, colour-vision-accessible palette,
adjustable text size, logical headings, keyboard-operable controls, and
tabular or plain-text alternatives to charts. The initial Streamlit
visualization prototype supplied by the team provides a useful starting
point for confidence tiers, accessible transcript rendering, summary
statistics, Altair charts, and downloadable tables. It will be integrated
selectively after its input schema is aligned with the application's
canonical transcript format. Audio playback, waveform interaction, and
direct pipeline integration are additional features to be developed.

\subsection{Application Component -- APP}

The application provides the user-facing workflow for processing and
reviewing recordings.

\begin{enumerate}
  \item \textbf{Upload:} select and upload a pre-recorded video or audio file.
  \item \textbf{Processing:} initiate analysis and view job status and
  stage-level progress.
  \item \textbf{Results:} open the generated transcript, speaker timeline,
  and confidence dashboard.
  \item \textbf{Review:} play selected regions, correct words, and revise
  speaker attribution where necessary.
  \item \textbf{Export:} download transcript results in structured JSON,
  plain text, and subtitle formats such as SRT or VTT.
  \item \textbf{History:} revisit previous recordings and their processing
  results.
\end{enumerate}

The proposed implementation uses Streamlit for the initial integrated
web interface, with Python services for media processing and persistence.
The interface and processing pipeline will communicate through clearly
defined functions and structured data schemas. If the final implementation
requires a separate API service, the same processing and data interfaces
can be exposed through REST endpoints. The application is intended for
single-user or small-team course-project deployment; a distributed
production-scale service is not a project requirement.

\subsection{Development Component -- DEV}

The development component ensures that the project is maintainable,
testable, and reproducible.

\begin{itemize}
  \item \textbf{Version control:} Git and GitHub with feature branches,
  meaningful commits, and code review among team members.
  \item \textbf{Modular design:} separate ingestion, preprocessing, VAD,
  ASR, alignment, diarization, fusion, confidence estimation, storage,
  and visualization modules.
  \item \textbf{Testing:} unit tests for data validation, timestamp
  handling, confidence calculations, transcript serialization, and
  evaluation metrics; integration tests for the end-to-end workflow.
  \item \textbf{Reproducible environment:} pinned or bounded Python
  dependencies, documented model versions, configuration files, and
  setup instructions. Docker may be used to simplify deployment.
  \item \textbf{Continuous integration:} GitHub Actions to run linting,
  static checks, and CPU-compatible tests on each relevant change.
  \item \textbf{Experiment tracking:} record model identifiers,
  configuration, dataset version, run time, and evaluation metrics.
  MLflow or a lightweight structured run log may be used.
  \item \textbf{Data and model versioning:} use versioned manifests and
  checksums for benchmark data and model configurations. DVC may be
  considered if dataset size and collaboration needs justify it.
  \item \textbf{Logging:} record stage status, duration, warnings, and
  failures without exposing uploaded audio or credentials in logs.
\end{itemize}

\section{Evaluation Plan}

Evaluation is divided into transcription quality, diarization quality,
confidence calibration, and application performance.

\subsection{Transcription and speaker metrics}

\begin{itemize}
  \item \textbf{Word Error Rate (WER):} measures substitutions,
  deletions, and insertions in the recognized transcript.
  \item \textbf{Character Error Rate (CER):} provides a complementary
  character-level measure where useful.
  \item \textbf{Diarization Error Rate (DER):} evaluates speaker
  segmentation and attribution when compatible reference RTTM annotations
  are available.
  \item \textbf{Speaker-attributed metrics:} cpWER and SA-WER assess
  transcription quality while accounting for speaker assignments.
\end{itemize}

\subsection{Confidence calibration}

Confidence quality will be evaluated by comparing predicted confidence
against observed word correctness. Planned measures include:
\begin{itemize}
  \item \textbf{Brier score:} mean squared error between predicted
  correctness probability and the binary correctness outcome.
  \item \textbf{Expected Calibration Error (ECE):} compares average
  confidence and empirical accuracy across confidence bins.
  \item \textbf{Reliability diagram:} plots empirical accuracy against
  predicted confidence.
  \item \textbf{Risk--coverage analysis:} measures error among the words
  retained when progressively excluding the least confident predictions.
\end{itemize}

All calibration results will state the evaluation split and number of
recordings/words used. Since the current labelled collection is small,
recording-level separation or cross-validation will be used to reduce
leakage, and calibration results will be interpreted with this limitation
in mind.

\subsection{System evaluation}

We will also measure processing time per minute of media, stage-wise
latency, memory/GPU usage where available, successful processing rate,
and the usability of the review and export workflow. Baselines will include
the ASR system without the proposed confidence calibration and, where
feasible, alternative model configurations.

\section{Expected Deliverables}

The expected course-project deliverables are:
\begin{enumerate}
  \item A working web application for uploading and processing pre-recorded
  video/audio.
  \item A modular speech-processing pipeline producing word-level
  timestamps and speaker-attributed transcripts.
  \item Word- and segment-level confidence estimates with a documented
  calibration and evaluation procedure.
  \item An interactive, accessible dashboard with transcript, confidence,
  speaker, timeline, and evaluation views.
  \item Persistent recording and result metadata, transcript correction
  history, and JSON/TXT/SRT/VTT exports.
  \item A benchmark report covering WER, available diarization metrics,
  calibration measures, and processing performance.
  \item Source code, tests, environment/setup documentation, and a
  demonstration video.
\end{enumerate}

\section{Scope, Assumptions, and Limitations}

\begin{itemize}
  \item The project processes pre-recorded media; real-time streaming is
  excluded.
  \item Pretrained models are integrated and evaluated; training a large
  ASR or diarization model from scratch is outside the scope.
  \item Reference transcripts are available only for the benchmark
  recordings. Arbitrary uploads will not have ground-truth-based accuracy
  measurements.
  \item Calibrated confidence is meaningful only to the extent supported
  by labelled validation data and a leakage-free evaluation procedure.
  \item Word-level alignment and speaker metrics depend on the quality
  and format of reference annotations. At present, structured JSON and
  RTTM references are available for only two benchmark recordings.
  \item Performance depends on hardware, model availability, and
  authentication requirements for gated pretrained models.
  \item Uploaded recordings may contain sensitive information. The
  application will document data handling, avoid logging media contents,
  and provide a clear retention policy for stored uploads.
  \item Accent, language, recording-device, and acoustic-condition
  differences may affect recognition and confidence reliability; these
  limitations will be discussed in the final evaluation.
\end{itemize}

\section{Conclusion}

Video Transcription with Confidence Analytics combines a practical
speech-intelligence application with the data-engineering and machine-learning
workflows required by CS698B. Its central contribution is not merely to
generate a transcript, but to expose where the transcript is uncertain,
measure whether confidence estimates are reliable, and help users review
the parts that need attention. The project brings together ETLS, pretrained
speech models, confidence calibration, accessible visualization, application
development, and reproducible engineering in one integrated system.

\end{document}
